% Part: computability
% Chapter: computability-theory
% Section: k-1

\documentclass[../../../include/open-logic-section]{subfiles}

\begin{document}

\olfileid{cmp}{thy}{k1}
\olsection{ન્યૂનીકરણક્ષમતાનું એક ઉદાહરણ}

\olref[ppr]{prop:reduce}નો એક ઉપયોગ જોઈએ.

\begin{prop}
\ollabel{prop:k1}
માનો કે
\[
K_1 = \Setabs{e}{\cfind{e}(0) \fdefined}.
\]
તો $K_1$ સંગણકીય રીતે પરિગણનીય છે, પરંતુ સંગણનીય નથી.
\end{prop}

\begin{proof}
$K_1 = \Setabs{e}{\lexists[s][T(e,0,s)]}$ હોવાથી,
\olref[eqc]{thm:exists-char} મુજબ $K_1$ સંગણકીય રીતે
પરિગણનીય છે.

$K_1$ સંગણનીય નથી તે બતાવવા, $K_0$ને તેમાં ઘટાડી શકાય
છે એવું બતાવીશું.

\begin{explain}
આ થોડું મુશ્કેલ છે, કારણ કે $K_1$ વડે આપણે ફક્ત એક
ચોક્કસ ઇનપુટ $0$થી શરૂ થતી સંગણનાઓ વિશે જ પૂછી શકીએ છીએ.
માનો કે આ પ્રકારના પ્રશ્નોના જવાબ આપી શકે એવો તમારો હોશિયાર
મિત્ર છે (આવા મિત્રોને ``ઓરેકલ'' કહે છે). પછી માનો કે
કોઈ તમને $\tuple{e, x}$ એ $K_0$માં છે કે નહીં,
એટલે મશીન $e$ ઇનપુટ $x$ પર થંભે છે કે નહીં, એવો પ્રશ્ન
પૂછે છે. તમે બીજું મશીન $e_x$ બનાવી શકો છો: તે
\emph{કોઈ પણ} ઇનપુટ લે, તેને અવગણે અને તેના બદલે
ઇનપુટ $x$ પર~$e$ ચલાવે. તેથી મશીન $e$ ઇનપુટ $x$ પર
થંભે છે કે નહીં તે પ્રશ્ન, મશીન $e_x$ ઇનપુટ $0$ પર
(અથવા બીજા કોઈ પણ ઇનપુટ પર) થંભે છે કે નહીં તે પ્રશ્નને
સમકક્ષ છે. એટલે તમે મિત્રને પૂછો કે આ નવું મશીન $e_x$
ઇનપુટ $0$ પર થંભે છે કે નહીં; સુધારેલા પ્રશ્નનો જવાબ
મૂળ પ્રશ્નનો જવાબ આપે છે. આમ $K_0$નું~$K_1$માં ઇચ્છિત
ન્યૂનીકરણ મળે છે.
\end{explain}

સર્વવ્યાપી આંશિક સંગણનીય વિધેયનો ઉપયોગ કરીને, $f$ આ રીતે
વ્યાખ્યાયિત ત્રિ-આર્ગ્યુમેન્ટવાળું વિધેય હોય:
\[
f(x,y,z) \simeq \cfind{x}(y).
\]
નોંધો કે $f$ તેના ત્રીજા ઇનપુટને સંપૂર્ણપણે અવગણે છે.
એવો સૂચક $e$ લો કે $f = \cfind{e}[3]$; એટલે
\[
\cfind{e}[3](x,y,z) \simeq \cfind{x}(y).
\]
$s$-$m$-$n$ પ્રમેય મુજબ એવું વિધેય $s(e,x,y)$ છે કે
દરેક $z$ માટે
\begin{align*}
\cfind{s(e,x,y)}(z) & \simeq \cfind{e}[3](x,y,z) \\
& \simeq \cfind{x}(y).
\end{align*}

\begin{explain}
ઉપરની અનૌપચારિક દલીલની ભાષામાં $s(e,x,y)$ એ એવા મશીનનો
સૂચક છે જે કોઈ પણ ઇનપુટ $z$ને અવગણીને $\cfind{x}(y)$
ગણે છે.
\end{explain}

ખાસ કરીને,
\[
\cfind{s(e,x,y)}(0) \fdefined \quad \text{ત્યારે અને માત્ર ત્યારે} \quad
\cfind{x}(y) \fdefined.
\]
બીજા શબ્દોમાં, $\tuple{x, y} \in K_0$ ત્યારે અને માત્ર
ત્યારે જ્યારે $s(e,x,y) \in K_1$. તેથી નીચે પ્રમાણે
વ્યાખ્યાયિત વિધેય $g$,
\[
g(w) = s(e,(w)_0,(w)_1)
\]
$K_0$નું~$K_1$માં ન્યૂનીકરણ છે.
\end{proof}

\end{document}
